\textbf{Phần 1. Trắc nghiệm lựa chọn}
\setcounterpageref{ex}{0}

\begin{ex} %[4D1N06]
Số gồm $5$ chục nghìn, $8$ nghìn và $4$ chục viết là
\choice
{\True $58\,040$}
{$58\,400$}
{$58\,004$}
{$50\,840$}
\loigiai{
Số đó là $58\,040$.\\
\textbf{Đáp án của câu hỏi là A}
}
\end{ex}

\begin{ex} %[4D1N06]
Trong các số: $124; 235; 346; 457$, số lẻ là
\choice
{$124$ và $346$}
{\True $235$ và $457$}
{$124$ và $235$}
{$346$ và $457$}
\loigiai{
$235$ và $457$ có chữ số tận cùng là 5 và 7 nên là số lẻ.\\
\textbf{Đáp án của câu hỏi là B}
}
\end{ex}

\begin{ex} %[4D1H06]
Giá trị của biểu thức $a + b \times 2$ với $a = 30, b = 15$ là
\choice
{\True $60$}
{$90$}
{$75$}
{$45$}
\loigiai{
$30 + 15 \times 2 = 30 + 30 = 60$.\\
\textbf{Đáp án của câu hỏi là A}
}
\end{ex}

\begin{ex} %[4D1H06]
Làm tròn số $68\,745$ đến hàng nghìn ta được
\choice
{$68\,000$}
{$68\,700$}
{\True $69\,000$}
{$70\,000$}
\loigiai{
Hàng trăm là 7 nên làm tròn lên $69\,000$.\\
\textbf{Đáp án của câu hỏi là C}
}
\end{ex}

\begin{ex} %[4D1H06]
Tìm $x$ biết: $x \times 4 = 48\,000$.
\choice
{$11\,000$}
{\True $12\,000$}
{$14\,000$}
{$16\,000$}
\loigiai{
$x = 48\,000 : 4 = 12\,000$.\\
\textbf{Đáp án của câu hỏi là B}
}
\end{ex}

\begin{ex} %[4D1H06]
Một hình vuông có cạnh $a = 12\text{ cm}$. Chu vi của hình vuông đó là
\choice
{$24\text{ cm}$}
{$144\text{ cm}$}
{$36\text{ cm}$}
{\True $48\text{ cm}$}
\loigiai{
$P = 12 \times 4 = 48\text{ cm}$.\\
\textbf{Đáp án của câu hỏi là D}
}
\end{ex}

\begin{ex} %[4D1V06]
Một cửa hàng ngày thứ nhất bán được $1250\text{ kg}$ đường, ngày thứ hai bán gấp đôi ngày thứ nhất, ngày thứ ba bán ít hơn ngày thứ hai $300\text{ kg}$. Cả ba ngày bán được số ki-lô-gam đường là
\choice
{\True $5950\text{ kg}$}
{$5800\text{ kg}$}
{$6000\text{ kg}$}
{$6200\text{ kg}$}
\loigiai{
Ngày 2: $1250 \times 2 = 2500\text{ kg}$.\\
Ngày 3: $2500 - 300 = 2200\text{ kg}$.\\
Tổng 3 ngày: $1250 + 2500 + 2200 = 5950\text{ kg}$.\\
\textbf{Đáp án của câu hỏi là A}
}
\end{ex}

\begin{ex} %[4D1V06]
Tính giá trị của biểu thức: $(a + b) \times c$ với $a = 125, b = 75, c = 8$.
\choice
{$1500$}
{$1700$}
{\True $1600$}
{$1800$}
\loigiai{
$(125 + 75) \times 8 = 200 \times 8 = 1600$.\\
\textbf{Đáp án của câu hỏi là C}
}
\end{ex}

\begin{ex} %[4D1V06]
Tổng của hai số chẵn liên tiếp là $150$. Số bé là
\choice
{$72$}
{$76$}
{$70$}
{\True $74$}
\loigiai{
Số bé: $(150 - 2) : 2 = 74$.\\
\textbf{Đáp án của câu hỏi là D}
}
\end{ex}

\begin{ex} %[4D1V06]
Một cuộn dây điện dài $100\text{ m}$. Người ta cắt lấy $4$ đoạn, mỗi đoạn dài $15\text{ m}$. Chiều dài đoạn dây điện còn lại là
\choice
{$35\text{ m}$}
{\True $40\text{ m}$}
{$45\text{ m}$}
{$50\text{ m}$}
\loigiai{
4 đoạn dài: $4 \times 15 = 60\text{ m}$.\\
Còn lại: $100 - 60 = 40\text{ m}$.\\
\textbf{Đáp án của câu hỏi là B}
}
\end{ex}

\begin{ex} %[4D1V06]
Một xí nghiệp may trong $4$ ngày may được $1600$ bộ quần áo. Hỏi trong $7$ ngày xí nghiệp đó may được bao nhiêu bộ quần áo?
\choice
{$2400$ bộ}
{$3000$ bộ}
{$3200$ bộ}
{\True $2800$ bộ}
\loigiai{
1 ngày: $1600 : 4 = 400$ bộ.\\
7 ngày: $400 \times 7 = 2800$ bộ.\\
\textbf{Đáp án của câu hỏi là D}
}
\end{ex}

\begin{ex} %[4D1C06]
Tìm số tự nhiên $x$ có hai chữ số biết rằng $x$ là số chẵn chia hết cho 5 và tổng hai chữ số của $x$ bằng 7.
\choice
{$25$}
{$52$}
{\True $70$}
{$61$}
\loigiai{
$x$ là số chẵn chia hết cho 5 nên chữ số tận cùng là 0. Do tổng hai chữ số bằng 7 nên chữ số hàng chục là 7. Số đó là $70$.\\
\textbf{Đáp án của câu hỏi là C}
}
\end{ex}

\textbf{Phần 2. Trắc nghiệm đúng sai}
\setcounterpageref{ex}{0}

\begin{ex} %[4D1H06]
Xét tính đúng sai của các biểu thức sau:
\choiceTF
{\True Số lớn nhất có 5 chữ số là $99\,999$.}
{Số liền trước của $50\,000$ là $49\,990$.}
{\True Số liền trước của $50\,000$ là $49\,999$.}
{\True Tổng của 3 số chẵn bất kì luôn là một số chẵn.}
\loigiai{
\begin{itemize}
    \item $99\,999$ là số lớn nhất có 5 chữ số. Đúng.
    \item Số liền trước của $50\,000$ là $49\,999 \neq 49\,990$. Ý 2 Sai, Ý 3 Đúng.
    \item Chẵn $+$ Chẵn $+$ Chẵn $=$ Chẵn. Đúng.
\end{itemize}
\begin{itemize}
    \item a) \textbf{Đúng}.
    \item b) \textbf{Sai}.
    \item c) \textbf{Đúng}.
    \item d) \textbf{Đúng}.
\end{itemize}
}
\end{ex}

\begin{ex} %[4D1V06]
Cho biểu thức $P = a \times 4 + b \times 4$. Xét tính đúng sai:
\choiceTF
{\True $P = (a + b) \times 4$.}
{\True Với $a = 15, b = 25$ thì $P = 160$.}
{Với $a = 10, b = 20$ thì $P = 100$.}
{\True Với $a = 10, b = 20$ thì $P = 120$.}
\loigiai{
\begin{itemize}
    \item Đặt thừa số chung: $P = (a + b) \times 4$. Đúng.
    \item Với $a=15, b=25 \Rightarrow P = 40 \times 4 = 160$. Đúng.
    \item Với $a=10, b=20 \Rightarrow P = 30 \times 4 = 120 \neq 100$. Ý 3 Sai, Ý 4 Đúng.
\end{itemize}
\begin{itemize}
    \item a) \textbf{Đúng}.
    \item b) \textbf{Đúng}.
    \item c) \textbf{Sai}.
    \item d) \textbf{Đúng}.
\end{itemize}
}
\end{ex}

\textbf{Phần 3. Trả lời ngắn}
\setcounterpageref{ex}{0}

\begin{ex} %[4D1H06]
Tìm số chẵn bé nhất có 5 chữ số khác nhau.
\shortans{10234}
\loigiai{
\textbf{Đáp số:} $10234$\\
Số đó là $10\,234$.
}
\end{ex}

\begin{ex} %[4D1H06]
Tính giá trị của biểu thức $a \times 6 - 45$ với $a = 15$.
\shortans{45}
\loigiai{
\textbf{Đáp số:} $45$\\
$15 \times 6 - 45 = 90 - 45 = 45$.
}
\end{ex}

\begin{ex} %[4D1V06]
Làm tròn số $85\,490$ đến hàng nghìn.
\shortans{85000}
\loigiai{
\textbf{Đáp số:} $85000$\\
Hàng trăm là 4 nên làm tròn xuống $85\,000$.
}
\end{ex}

\begin{ex} %[4D1V06]
Hai số lẻ liên tiếp có tổng bằng $84$. Tìm số lớn.
\shortans{43}
\loigiai{
\textbf{Đáp số:} $43$\\
$(84 + 2) : 2 = 43$.
}
\end{ex}

\begin{ex} %[4D1V06]
Có $4$ thùng nước mắm, mỗi thùng chứa $25\text{ l}$. Người ta đóng đều số nước mắm đó vào các chai $2\text{ l}$. Hỏi đóng được bao nhiêu chai nước mắm?
\shortans{50}
\loigiai{
\textbf{Đáp số:} $50$\\
$4 \times 25 : 2 = 100 : 2 = 50$ chai.
}
\end{ex}

\begin{ex} %[4D1C06]
Tìm giá trị của biểu thức $A = 100 - a \times b$ biết $a$ là số lẻ nhỏ nhất có hai chữ số khác nhau và $b$ là số chẵn lớn nhất có một chữ số.
\shortans{4}
\loigiai{
\textbf{Đáp số:} $4$\\
$a = 13$ (số lẻ bé nhất có 2 chữ số khác nhau), $b = 8$.\\
$A = 100 - 13 \times 8 = 100 - 104$ (nếu $a = 11 \Rightarrow 100 - 88 = 12$; nếu $a = 12$ chẵn; nếu $a = 13, b = 6 \Rightarrow 100 - 78 = 22$; nếu $a = 12, b = 8 \Rightarrow 100 - 96 = 4$).
}
\end{ex}
